\subsection{Atomic input logic and possible winners}
\label{app:atomic-inputs}

Atomic valuations restrict the inputs, while intermediate values can still
be ambiguous. This distinction determines the appropriate proof system.

\begin{proposition}[Atomic validity and substitution]
\label{olv:prop:atomic-substitution}
For $K\geq3$, the equations valid at every atomic valuation are not the
consequences of any ordinary one-sorted equational theory, even an infinite
one. Substitution by atom-preserving terms is admissible.
\end{proposition}
\begin{proof}
The equation $x\del y=y$ holds whenever $y$ is atomic. Substitute
$u\star_c v$ for $y$, and assign $x=\{c\}$, $u=\{a\}$, $v=\{b\}$,
for distinct $a,b,c$. The substituted right side is $\{a,b\}$, whereas
the left side is $\{a,b,c\}$. Ordinary equational consequence is closed
under this substitution, so it cannot coincide with atomic validity.
An atom-preserving substitution, in contrast, sends every atomic input
valuation to an atomic valuation of the original variables, proving
admissibility directly.
\end{proof}

There is a finite equational presentation in a sorted case expansion.
Use a role sort $A=I$ for input variables and a winner sort $W=\Pow$
for intermediate values. Write $\iota:A\to W$ for singleton injection,
name every role and every winner value $v_S$, and add
\[
 \operatorname{case}(x;(w_r)_{r\in I})=w_x.
\]
For fixed $K$ the following are finitely many equations, with all
substitutions respecting the two sorts:
\begin{align*}
 \operatorname{case}(r;(w_i)_i)&=w_r,&
 \operatorname{case}(x;(w)_i)&=w,\\
 \iota x&=\operatorname{case}(x;(v_{\{r\}})_r),&
 \iota r&=v_{\{r\}},\\
 \operatorname{case}(x;(u_i)_i)\star_r
   \operatorname{case}(x;(w_i)_i)
   &=\operatorname{case}(x;(u_i\star_r w_i)_i),\\
 \pi\bigl(\operatorname{case}(x;(w_i)_i)\bigr)
   &=\operatorname{case}(x;(\pi(w_i))_i).
\end{align*}
Include also the ground operation tables on the named values $v_S$.

\begin{theorem}[A finite sorted presentation]
\label{olv:thm:atomic-calculus}
For native terms $s,t$, the equation $s=t$ holds at every atomic
valuation if and only if these finite sorted equations derive the
translations of $s,t$ that replace each input $x$ by $\iota x$.
Substitution by any atom-preserving native terms is admissible in this
presentation.
\end{theorem}
\begin{proof}
Every displayed equation is sound in the indicated two-sorted algebra.
For completeness, the injection and diagonal equations and the operation
distribution equations derive the case expansion of a native term at
any input variable. Expand all its input variables in a fixed order.
The leaves are ground terms, which the finite tables reduce to their
named winner values. Two atom-valid terms therefore reduce to the same
ordered case tree. Admissibility follows from the preceding proposition
and completeness.
\end{proof}

This theorem presents the native atomic equations in the sorted expansion;
it does not assert a finite basis for the winner algebra with unrestricted
winner-sort inputs.

The possible-win question has a different answer depending on whether
internal role renamings are allowed.

\begin{theorem}[Atomic possible-win complexity]
\label{olv:thm:atomic-sat}
Fix $K\geq2$ and a queried role $a$.
\begin{enumerate}[label=(\roman*)]
\item For terms without internal role renamings, some atomic valuation
makes $a$ a possible winner if and only if the valuation assigning
$\{a\}$ to every variable does. This gives a linear-time test at fixed $K$.
\item For the full signature, deciding whether some atomic valuation
makes $a$ a possible winner is NP-complete.
\end{enumerate}
\end{theorem}
\begin{proof}
For (i), use simultaneous induction on a public term: replacing all its
inputs by $\{a\}$ preserves possible membership of $a$, and if the new
value is a foreign singleton $\{r\}$, the old value was already that
singleton. At an $a$-choice, membership of $a$ is the disjunction of its
membership in the children. At a foreign $r$-choice, an absorbing
$\{r\}$ after replacement was already absorbing before replacement;
otherwise the children take union. The singleton clause follows from
the corresponding disjunction for the owning role and conjunction for
every other role. These facts prove the simultaneous induction.

For (ii), membership in NP follows by guessing one role per occurring
variable and evaluating the term. Choose $b\ne a$. There is a fixed
unary term $q$ with
\[
 q(\{r\})=\begin{cases}\{a\},&r=a,\\\{b\},&r\ne a.\end{cases}
\]
Explicitly, fold $\star_b$ over $x$ and $(b\ c)(x)$ for all
$c\notin\{a,b\}$. An input $a$ stays $a$; every other atomic input
has a copy equal to $b$, which absorbs the fold.
Translate a Boolean CNF by replacing a positive literal $x$ by $q(x)$,
a negative literal by $(a\ b)(q(x))$, conjunction by $\star_b$, and
disjunction by $\star_a$. Empty conjunctions and disjunctions use the
winner constants $a,b$. The resulting term always returns one of these
two atoms on atomic inputs, and returns $\{a\}$ precisely when the
Boolean assignment $x\mapsto[x=a]$ satisfies the CNF. Every Boolean
assignment has such an atomic encoding. The translation has linear size
at fixed $K$, proving NP-hardness.
\end{proof}

